\documentclass[11pt]{article}
\usepackage[margin=1in]{geometry}
\usepackage{amsmath,amssymb,amsfonts,amsthm,mathtools,bm}
\usepackage[T1]{fontenc}
\usepackage{microtype}
\newtheorem{theorem}{Theorem}
\newtheorem{lemma}{Lemma}
\newtheorem{cor}{Corollary}
\newtheorem{defn}{Definition}
\newtheorem{principle}{Principle}
\newtheorem{statement}{Statement}
\theoremstyle{remark}
\newtheorem{statementr}{Statement}

\title{A Very Useful Frame Field}
\author{James Taylor}
\date{Working draft derived from the 2003 thesis}
\begin{document}
\maketitle

\section{A very useful frame field} \label{a:difgeom} 


In this part, we shall explain how to find sections satisfying
constraints on the symmetric part of their covariant derivative. 
Let $M$ be a manifold with connection and $E$ be a vector bundle over $M$
whose fiber is either isomorphic to \ensuremath{\mathbb{R}^{k}} or $\ensuremath{\mathbb{C}}^k$. Fix $p
\in M$. This appendix relies upon the work of Chapter \ref{c:taylor}
and is used in Chapter \ref{c:geom}


Choose a coordinate neighborhood about $p$ such that for every point
$q$ in the neighborhood, there is a unique
geodesic between $p$ and $q$. Choose a basis $\{e_i\}_{i=1}^{k}$
of the fiber $E_p$ for the rank $k$ bundle $E$  
over $M$. Parallel transport these vectors along the geodesics
emanating from $p$, defining a basis field  also denoted by 
$\{e_i\}_{i=1}^{k}$. 

Let $q$ be any point in the neighborhood. Take
\ensuremath{\gamma_v}\ to be the unique geodesic between $p$ and $q$. Then, by
definition, for each $i$, $\nabla_v [e_i] = 0$ along $\ensuremath{\gamma_v}$. Thus,
Corollary~{\ref{tc:da}} applies and we may conclude that the {\bf symmetric}
 part of  $D^n_p [e_i]
=0$ for all $i$.\footnote{It is true that the diagonal part of a
  tensor determines its symmetric values. One proof is based on the
  following. Define $f(a_1, \ldots,a_k) \ensuremath{:=} A(\{\sum_i a_i e_i \}^n)$.
Notice that we know $f$ for all values of the $a_i$'s. Expand the
right-hand side. Taking the appropriate \ensuremath{\text{n}^{\text{th}}}-order partial derivatives
with respect to the 
$a_i$'s and evaluating the derivatives at $(0,\ldots, 0)$, we can
generate any symmetric expression in the $e_i$'s of length $n$ that we
wish.} 
 We emphasize that this is only true at $p$.

Now define $A := \sum_{i=1}^k f^i e_i$ where the $f^i$ are
scalar functions defined on the neighborhood. Then the symmetric part
of the covariant
derivatives of $A$ are the symmetric covariant derivatives of the $f^i$'s. 
Using Corollary~{\ref{tc:imp}}, we know that the symmetric part of the
higher order covariant derivatives of the $f^i$'s are  the
Euclidean derivatives when expressed in 
 normal coordinates. Thus, it is reasonably
easy to specify the symmetric part of various higher order covariant
derivatives. 

The other parts of the covariant derivatives are part of the geometric
structure of the connection. For example, the second covariant
derivative's antisymmetric part is the curvature tensor, at least if
the connection on $M$ has no torsion. More specifically, we claim
\[
\ensuremath{{{\nabla}}_{x}^{}{[\ensuremath{{{\nabla}}_{y}^{}{[A]}}{}]}} - \ensuremath{{{\nabla}}_{y}^{}{[\ensuremath{{{\nabla}}_{x}^{}{[A]}}{}]}}
- \ensuremath{{{\nabla}}_{[x,y]}^{}{[A]}} 
=
D^2[A] (y\otimes x -x \otimes y) - D[A] (\tau(x, y))
\]
where $\tau$ is the torsion; torsion, acting on a function $f$, is 
\[
\tau(x,y)f = (\ensuremath{{{\nabla}}_{x}^{}{[y]}} - \ensuremath{{{\nabla}}_{y}^{}{[x]}} -[x,y])f 
= D^2[f](x\otimes y -y \otimes x).
\]
These assertions may be verified using the chain rule. A typical fact
to use, for example, is
\[
\ensuremath{{{\nabla}}_{x}^{}{[\ensuremath{{{\nabla}}_{y}^{}{[A]}}{}]}} =  
D^2 [A] (y \otimes x) + D[A](\ensuremath{{{\nabla}}_{x}^{}{[y]}}).
\]
The formulas involving the differentiation of the vector fields are
the usual formulas for the curvature and torsion of a
connection. Writing it in this fashion demonstrates that curvature is
independent of the connection on the manifold. But writing it using
the higher covariant derivatives allows us to immediately see that the
objects are tensors in the tangent vector slots involving $x$ and
$y$. As for the curvature being a tensor in $A$, that is why the
torsion term is necessary. Indeed, multiply $A$ by a function $f$, and
one sees that cross terms cancel when computing the antisymmetric part
leaving the expression in $A$ that we want plus a torsion
term. Subtracting the torsion therefore eliminates the contribution of
$f$. 

Another way to view the subtraction of the torsion
 is to think of torsion as a type of
curvature and that multiplication of a scalar function times a section
is actually the tensor product of the section of $E$ with a section of
the trivial one-dimensional bundle over $M$. Then, as is easily seen,
the curvature of the tensor product connection  is a sum of the two
 curvatures. Thus, we subtract the function's curvature to obtain the
 curvature of the bundle $E$. 

In any event, when we use our special basis, we are only left with the
antisymmetric part  as the symmetric part of
the basis fields vanish at $p$. Thus the second covariant derivatives of these
fields is the curvature tensor applied to field that we are
differentiating, assuming that there is no torsion. 

As an application of this, we can prove that if $E$ is a complex
bundle with a connection and a parallel inner product $\ensuremath{({}, {})}$,
then the curvature tensor is anti-self-adjoint. Let $R$ be the
curvature tensor. Then, to establish our claim, we need to show that 
\[
\ensuremath{({\ensuremath{\psi}}, {R\phi})} = - \ensuremath{({R \ensuremath{\psi}}, {\phi})}
\]
for any sections $\ensuremath{\psi}$ and $\phi$ of $E$. Fix a point $p \in M$. 
Given $\ensuremath{\psi}_p, \phi_p \in
E_p$, define $\ensuremath{\psi}$ and $ \phi$  by parallel transporting the value at
$p$ to all the other points in an appropriately small neighborhood
about $p$. As just explained, $D^2 \ensuremath{\psi} = R\ensuremath{\psi}$ and similarly for
$\phi$ where we use a torsionless connection on $M$ to compute the
second covariant derivative. Since the inner product is parallel, we
have that $\ensuremath{({\ensuremath{\psi}}, {\phi})}_q = \ensuremath{({\ensuremath{\psi}}, {\phi})}_p$ for all $q$ in the
neighborhood. More to the point, we have $D[\ensuremath{({\ensuremath{\psi}}, {\phi})}] =0$ and,
more importantly, 
\[
D^2 [\ensuremath{({\ensuremath{\psi}}, {\phi})}] = 0.
\]
But notice that 
\[
D^2_p[\ensuremath{({\ensuremath{\psi}}, {\phi})}] = \ensuremath{({D_p^2[\ensuremath{\psi}]}, {\phi})} + \ensuremath{({\ensuremath{\psi}}, {D_p^2[\phi]})}
= \ensuremath{({R \ensuremath{\psi}}, {\phi})}_p + \ensuremath{({\ensuremath{\psi}}, {R\phi})}_p;
\]
note that only at $p$ do we have the vanishing of the first derivative
terms in addition to the second derivative terms being the curvature tensor. 
As one can see, we have our result.

\clearpage
\bibliographystyle{abbrv}
\bibliography{references}
\end{document}
